\section{Cách viết tắt và mở rộng hệ thống lô-gích}

\subsection{Viết tắt các lập luận}

Bạn đọc đã có thể cảm nhận rằng xây dựng và giải thích cách hoạt động của ngôn ngữ hình thức trong các phần trước khá là dài dòng. Điều này là hiển nhiên, khi chúng ta đang chuẩn tắc hóa các lập luận thông thường. Nhưng, trong thực hành sau này, chúng ta sẽ ít khi chỉ rõ một cách tường minh chúng ta đang sử dụng những tính chất hay quy tắc lập luận diễn dịch nào. Mặc định rằng, chúng ta coi như bạn đọc sẽ có thể tự đưa ra cơ sở cho từng bước lập luận, cũng giống như chúng ta không nghiên cứu một bài văn hay bài thơ mà cứ đâm đầu vào từng chữ cái một vậy.

Hãy xem xét cách viết tắt của một số cấu trúc chứng minh quen thuộc. Khi muốn khẳng định \begin{equation*}
   \forall y, B(y) \implies C(y),
\end{equation*}
người ta thường bắt đầu với một trong những cách viết:
\begin{itemize}
   \item ``gọi $y$ sao cho $B(y)$'',
   \item ``xét $y$ thỏa mãn $B(y)$'',
   \item ``lấy $y$ để $B(y)$'',
\end{itemize}
hoặc nhiều cách khác. Một điều không hay lắm là khi có $\exists y, B(y)$, người ta cũng có thể dùng cách viết y hệt. Tác giả không có một phương pháp nào để có thể phân biệt ý của người viết. Bạn đọc cần phải dựa vào ngữ cảnh để xem ý tưởng đằng sau những câu như vậy.

\subsection{Mở rộng các lập luận}

Có thể sẽ có câu hỏi rằng liệu có thể tối giản lại ngôn ngữ hình thức mà chúng ta vừa xây dựng. Bạn đọc mà tìm hiểu thì sẽ thấy rằng ngôn ngữ được xây dựng thuộc về giải tích vị từ bậc hai. Nếu bỏ đi khả năng kết hợp lượng từ với vị từ hay hàm, chúng ta sẽ có vị từ bậc nhất. Ngoài ra, khi nhận ra rằng có hai bậc lô-gích, một số nghi vấn có thể được đặt ra về vấn đề sự tồn tại của lô-gích trên các bậc cao hơn. Câu trả lời cho câu hỏi đó là có, nhưng khó tưởng tượng. Bạn đọc muốn tìm hiểu có thể xem sự xây dụng lô-gích bậc cao của Chớc\footnote{Alonzo Church (1903 - 1995).}.

Hơn thế nữa, để ý rằng, chúng ta mới xây dựng các quy tắc và phương pháp lập luận trên một phạm vi toán học rộng nhất có thể. Chúng ta chưa đề cập tới bất kì tính chất nào của số hay hình. Chúng ta còn chưa có định nghĩa để khẳng định $$\forall x, x = x.$$

Cứ đi dần vào sâu toán học, hay thế giới vật lí, chúng ta sẽ nhận ra rằng không có một bộ tiên đề, quy luật hay chỉ dẫn nào có thể bao quát hết tất cả các tính chất. Tuy nhiên, trong thực tế, nhiệm vụ của chúng ta không phải là đi tìm cách chuẩn tắc hóa tất cả mọi thứ. Đó không phải là phạm trù của khoa học. Hãy nhớ rằng, chúng ta mới chỉ đi qua lập luận diễn dịch, còn cần phải xem xét làm sao để thực hiện lập luận quy nạp\footnote{Bạn đọc sẽ nhận ra rằng lập luận quy nạp sẽ phức tạp hơn lập luận diễn dịch nhiều, do nó cần phải dựa vào trải nghiệm cá nhân mà những cái đó có xác suất không phản ánh đúng thực tiễn (báo hiệu trước đến công cụ chúng ta cần phải xây dựng). Đấy cũng có thể là một lí do mà bạn đọc có thể biết đến ``thần đồng Toán học'' có thể giải được những bài toán khó (môn có tính tiền định cao), nhưng chắc chắn là không bao giờ có ``thần đồng Sinh học'' tự nhiên phát hiện được quy luật di truyền hay ``thần đồng Địa lí học'' giải thích được sự chuyển dịch của các dân tộc (những môn cần có khảo sát thực nghiệm kiểm chứng).}. Do đó, việc của chúng ta bây giờ, giống như bổ sung từ và cách sử dụng vào trong ngôn ngữ, để mở rộng các phương pháp lập luận phục vụ các mục đích sau này, là xây dựng các cấu trúc toán học mới và bổ sung thêm kiến thức thực tiễn.